% Laget automatisk av elevark/lag_elevark.py fra opplegg/ballstabel.html. Endringer her blir overskrevet.
\documentclass[11pt]{article}
\usepackage{elevark}
\begin{document}
\elevarkhode{Fysikk 1 · Bevegelsesmengde og støt}{Ballstabelen}{https://fysikklab.no/opplegg/ballstabel.html}{Opplegget på nett}

Hold en tennisball oppå en basketball og slipp begge samtidig. Tennisballen skyter til værs, mye høyere enn den ble sluppet fra, mens basketballen nesten blir liggende igjen. Hvor mange ganger høyere spretter tennisballen, og bryter det med energibevaringen?

\section*{Hensikt}

\begin{itemize}
\item Måle hvor mye høyere den øverste ballen spretter, og sammenligne med en enkel modell.
\item Føre regnskap over den mekaniske energien før og etter.
\end{itemize}

\section*{Utstyr}

\begin{itemize}
\item En basketball (eller fotball) og en tennisball (eller sprettball), og en vekt.
\item Målebånd teipet loddrett på en vegg, og en mobil som kan filme i sakte film.
\end{itemize}

\section*{Sikkerhet}

Tennisballen kan sprette flere meter opp og ut til siden. Gjør forsøket der det er høyt under taket, unna lamper og vinduer, og ikke len dere over ballene.

\section*{Framgangsmåte}

\begin{enumerate}
\item Vei ballene.
\item \textbf{Hver ball alene:} Slipp hver ball fra 70~cm og mål hvor høyt den spretter. Mål fra underkanten av ballen.
\item \textbf{Ballstabelen:} Legg basketballen på gulvet og merk av høyden til toppen av den. Den høyden er nullpunktet for tennisballen.
\item Hold tennisballen oppå basketballen og løft begge slik at underkanten av basketballen er i høyden $h_{0}$. Slipp samtidig, så rett som mulig.
\item Film mot målebåndet. Les av hvor høyt tennisballen kom over nullpunktet $(h_{\mathrm{t}})$ og hvor høyt basketballen spratt $(h_{\mathrm{b}})$.
\item Gjenta for fem ulike slipphøyder. Ta flere slipp per høyde og bruk det beste, for stabelen bommer lett.
\end{enumerate}

\section*{Oppgaver}

\begin{oppgaveliste}
\oppgave \fase{gjett}{Gjett}Stabelen slippes fra 50~cm. Hvor høyt tror dere tennisballen spretter? Hvor høyt spretter basketballen?
\svar
\oppgave \fase{test}{Mål}Slipp ballene hver for seg fra 70~cm. Hvor stor del av energien er igjen etter ett sprett for hver av dem?
\svar
\oppgave \fase{test}{Mål}Gjør målingene med stabelen, eller trykk på «Fyll inn eksempel». Hvor mange ganger høyere $enn h_{0}$ spretter tennisballen?
\svar
\oppgave \fase{forklar}{Forklar}Tenk deg at basketballen er mye tyngre enn tennisballen og at alle støt er elastiske. Basketballen har snudd og går oppover med farten $v$ når den møter tennisballen, som går nedover med farten $v$. Bruk at tennisballen spretter fra basketballen med samme relative fart som den kom inn med, og vis at tennisballen får farten $3v$. Hvor høyt kommer den da?
\svar
\oppgave \fase{forklar}{Vurder}Tennisballen kommer høyere enn der den startet. Bryter det med bevaring av energi? Hvor mye av energien er igjen i målingene deres, og hvor har resten tatt veien?
\svar
\oppgave \fase{utfordring}{Utfordring}Med tre baller, hver mye lettere enn den under, og bare elastiske støt: Hvor fort går den øverste ballen, og hvor mange ganger høyere $enn h_{0}$ kan den komme?
\svar
\end{oppgaveliste}

\end{document}
